% =============================================================================
% SECTION III — Factor Selection and Spanning Regressions
% =============================================================================


% ── SEMAPHORE ──────────────────────────────────────────────────────────────────

The benchmark factor models of Section~\ref{sec:benchmarks} are
designed to capture known systematic exposures, but they leave open
the question of whether a broader set of risk factors can better
explain strategy excess returns and absorb the residual alpha.
Starting from the candidate universe of $75$ tradable risk factors
described in Section~\ref{sec:factor_universe}, we proceed in two
steps. Section~\ref{sec:pca} extracts a small number of latent factors
by principal components, following \citet{connor1986performance,
connor1988risk}, and tests whether these diffuse combinations span
strategy excess returns. Section~\ref{sec:aen} turns to a supervised
alternative---best-subset selection \citep{bertsimas2016best}---that
picks the individual factors most relevant for each strategy from the
same candidate set. As in Section~\ref{sec:benchmarks}, each layer is
subjected to the conditional alpha model
(Section~\ref{sec:conditional_alpha}) and to an out-of-sample
re-estimation of the hedge, and
Section~\ref{sec:factor_discussion} consolidates the alpha estimates
across all layers of the paper.


% ── III.A CANDIDATE FACTOR UNIVERSE ───────────────────────────────────────────

\subsection{Candidate Factor Universe}
\label{sec:factor_universe}

We assemble a candidate set of $75$ risk factors from a comprehensive
review of the empirical asset pricing, intermediary, and macro-finance
literatures. Each factor is selected on the basis of a published
theoretical or empirical link to fixed-income arbitrage returns,
intermediary balance-sheet constraints, or broad market risk
conditions. The factors are organised into five economic categories:
credit risk (Panel~A), liquidity and funding (Panel~B), equity
(Panel~C), volatility (Panel~D), and interest rates and term
structure (Panel~E). The complete list with
definitions, data sources, and references is reported in
Table~\ref{tab:factor_list} (Appendix~\ref{app:factor_list}). This
candidate universe is the common input to both the PCA of
Section~\ref{sec:pca} and the best-subset selection of
Section~\ref{sec:aen}.

Every candidate is tradable: each series is the excess return on an
implementable position, such as a cash index, a long-short factor
portfolio, a derivative overlay, or a DV01-neutral curve trade. State variables
enter through traded counterparts rather than as raw proxies; the
intermediary capital factor of \citet{he2017intermediary} enters
through their tradable primary-dealer equity factor, which the authors
show tracks innovations in the intermediary capital ratio with a
correlation of $0.92$, and aggregate liquidity enters through the
tradable liquidity factor of \citet{pastor2003liquidity}. As in
Section~\ref{sec:benchmarks}, dollar-denominated factors---including
these two---enter as euro excess returns, converted following
\citet{gluck2021currency}; factors with native euro-area underlyings
enter directly. As in Section~\ref{sec:benchmarks}, every regressor is a realised
excess return, so each intercept $\alpha_i$ retains its interpretation
as a Jensen alpha.


% ── III.B PRINCIPAL COMPONENT ANALYSIS ────────────────────────────────────────

\subsection{Principal Component Analysis}
\label{sec:pca}

When the number of candidate factors is large relative to the sample
size, regressing strategy returns on the full set is unreliable: the
strong collinearity of the candidates inflates the variance of the
coefficient estimates and overfits the sample. The standard response
is to summarise the common variation in a small number of principal
components and test whether these diffuse combinations span strategy
returns; estimating latent factors from a large panel and using them
to measure the abnormal return of a managed position is the
performance-measurement framework of \citet{connor1986performance,
connor1988risk}.

We estimate the spanning regression in two forms. In sample, the
components are estimated once on the full factor panel and used as
observable factors; with the cross-sectional dimension $N$ and the
time dimension $T$ both large, this yields consistent second-stage
estimates of $\alpha_i$ and $\boldsymbol{\beta}_i$
\citep{giglio2021asset}. Out of sample, reported below, the components
and the betas are re-estimated on a past-only window so that the hedge
uses no future information; this is the test of whether the alpha
would have survived in real time. We
standardise each factor to zero mean and unit variance over the
estimation sample and retain the first $K = 8$ components, the number
selected by the \citet{bai2002determining} $IC_{p2}$ criterion (the
less parsimonious $IC_{p1}$ indicates thirteen). The eight components
capture $64.7\%$ of the variance in the balanced factor panel
(Figure~\ref{fig:pca_scree}), and robustness of the alpha to
$K \in \{5, 8, 11, 13\}$ is reported in Appendix~\ref{app:pca_wk}.

The spanning regression takes the form
\[
r_{i,t} \;=\; \alpha_i \;+\; \boldsymbol{\beta}_i'\,
\mathbf{PC}_t \;+\; \varepsilon_{i,t},
\]
where $r_{i,t}$ is the monthly excess return of strategy $i$ and
$\mathbf{PC}_t$ is the $K \times 1$ vector of contemporaneous
principal component scores. Inference uses Newey--West HAC standard
errors \citep{newey1987simple}.

Table~\ref{tab:pca_spanning} reports the estimates. The intercept
remains positive and significant at the $1\%$ level for all three
strategies: $+1.44\%$ p.a.\ for BTP~Italia, $+4.78\%$ for CDS--Bond
Basis, and $+2.21\%$ for iTraxx Skew. The adjusted $\bar{R}^2$ is
low throughout, reaching at most $13.4\%$: the diffuse principal
components do not span strategy excess returns.
Table~\ref{tab:pca_loadings} reports, for each retained component, the
candidate factors with the largest absolute loading. The components
admit a clear economic reading: PC1 is dominated by credit and
equity-market factors, PC2 by global rates and bond returns, and PC3
by swap spreads and dealer CDS.

\paragraph{Out-of-Sample Alpha.}

As in Section~\ref{sec:benchmarks}, the full-sample loadings could
reflect an in-sample-optimal combination of the factors---and, since
the components are themselves estimated, could carry a
generated-regressor bias---rather than a return an investor could have
earned in real time \citep{moreira2017volatility,
cederburg2020performance}. Panel~A of
Table~\ref{tab:pca_oos} addresses this directly. At each evaluation
date we re-estimate the principal components and the spanning
betas on a past-only window---expanding (initialised at $60$ months)
as the baseline and a $60$-month rolling window as a robustness
check---and define the out-of-sample alpha as the mean of the hedged
return, tested with the \citet{newey1994automatic} HAC correction.
Because the components and the hedge coefficients are estimated in the
same window, the out-of-sample alpha is invariant to the sign and
rotation indeterminacy of the principal components: any orthogonal
rotation of the basis cancels between the loadings and the scores. 
The out-of-sample alpha remains positive and
significant at the $1\%$ level for all three strategies: $+1.61\%$
p.a.\ for BTP~Italia, $+4.55\%$ for CDS--Bond Basis, and $+1.85\%$ for
iTraxx Skew under the expanding window, close to the in-sample
estimates of Table~\ref{tab:pca_spanning}, so the alpha is not an
artefact of full-sample factor estimation.

Because PCA is an unsupervised method---it extracts factors that
maximise variance in the candidate set, not factors that explain
variation in strategy returns---it cannot identify strategy-specific
risk drivers. This motivates the supervised selection of the next
subsection.


% ── III.C BEST-SUBSET SELECTION ───────────────────────────────────────────────

\subsection{Best-Subset Selection}
\label{sec:aen}

\paragraph{Preprocessing.}

The candidate factors are preprocessed for selection. Six candidates
are near-duplicates of another factor in the universe, with absolute
correlation above $0.95$; one member of each pair is removed, leaving
$p = 69$ candidates in the selection design. Each factor is
standardised to unit $L_2$-norm, so that which factors enter is
invariant to scale; point estimates and inference are subsequently
computed on the original units. Even at $p = 69$, the space of
possible factor subsets---$2^{69} \approx 6 \times 10^{20}$---rules
out exhaustive search and calls for a selection method with an
explicit cardinality control.

\paragraph{Selection.}

For each strategy we solve the best-subset problem of
\citet{bertsimas2016best},
\begin{equation}\label{eq:best_subset}
\min_{\alpha,\,\boldsymbol{\beta}}\;
\frac{1}{T}\sum_{t=1}^{T}\bigl(r_t - \alpha -
\boldsymbol{\beta}'\mathbf{f}_t\bigr)^2
\quad \text{subject to} \quad
\|\boldsymbol{\beta}\|_0 \le k,
\end{equation}
computed with the splicing algorithm of \citet{zhu2020polynomial}
polished against a greedy forward search, retaining whichever feasible
$k$-subset attains the lower in-sample residual sum of squares. The $\ell_0$ constraint has two properties
that penalised methods lack and that matter in our design. First, a
factor that enters is not shrunk, so among a cluster of correlated
candidates the constraint keeps the most informative member and drains
its near-duplicates, instead of splitting a common signal across the
cluster or shrinking it toward zero as the LASSO does. Second, the
cardinality is controlled directly: we fix $k = 8$ for parsimony and
cross-strategy comparability, rather than selecting the count by an
information criterion or cross-validation, whose objective is nearly
flat---and whose argmin is therefore ill-determined---on
alpha-dominated series with weak factor signal.
Table~\ref{tab:ml_ksens} (Appendix~\ref{app:aen_sensitivity})
re-estimates the selection at $k \in \{5, 8, 10\}$: the alpha stays
significant at the $1\%$ level and close to its baseline value at every
cardinality, so the conclusions do not hinge on the choice.

\paragraph{Post-Selection Inference.}

Given the selected set $\hat{S}_i$, we re-estimate the spanning
regression by unrestricted OLS on the original (un-standardised)
factors,
\begin{equation}\label{eq:post_ols}
r_{i,t} \;=\; \alpha_i \;+\; \boldsymbol{\beta}_i'\,
\mathbf{f}_{\hat{S}_i,t} \;+\; \varepsilon_{i,t},
\end{equation}
with Newey--West HAC standard errors. Because the set is selected
on the same data, we complement the analytic $t$-tests with a
stationary block bootstrap \citep{politis1994stationary} of the
post-selection regression ($9{,}999$ replications, expected block
length nine months) and report the percentile $95\%$ confidence
interval for alpha. Because selecting controls from the data can bias
the surviving alpha, we also report the post-double-selection control
set of \citet{belloni2014inference}, which selects controls both for
the strategy return and for the factors of interest; the set appears
alongside the selected sets in Appendix~\ref{app:aen_selection_matrix}
and is discussed in the robustness analysis below.

\paragraph{Results.}

Table~\ref{tab:aen_stable_ols} reports the post-selection estimates.
Alpha is significant at the $1\%$ level for all three strategies:
$+2.2\%$ p.a.\ for BTP~Italia (bootstrap CI $[1.3, 3.0]$), $+4.8\%$
for CDS--Bond Basis (CI $[2.9, 6.6]$), and $+2.9\%$ for iTraxx
Skew (CI $[1.5, 3.8]$). With eight factors each, the models
deliver adjusted $\bar{R}^2$ of $31\%$, $21\%$, and $25\%$---well
above any benchmark framework of Section~\ref{sec:benchmarks} and the
PCA above---yet absorb none of the three alphas.

The selected sets differ across strategies in ways that mirror each
trade's structure, and all but two of the twenty-four selected
loadings are significant. The BTP~Italia set maps onto the three ingredients of the trade:
inflation, liquidity, and sovereign credit. Long-horizon inflation
expectations carry one of the strongest loadings ($-0.23$,
$t = -2.7$)---the natural exposure of the retail inflation hedge
described in Section~\ref{sec:interdependencies}---and the
\citet{pastor2003liquidity} traded liquidity factor enters alongside,
consistent with a basis earned on an illiquid, buy-and-hold retail
instrument. The credit block combines the fixed-income defensive
factor ($-0.21$, $t = -4.9$), the euro default premium, and the
five-year swap-spread return, the exposures of a position long and
short the same sovereign; the equity market ($-0.04$, $t = -4.0$),
size, and fixed-income value complete the set. The CDS--Bond set
maps onto the credit and intermediary channels: the strategy loads
negatively on the prime-broker CDS factor ($-0.22$, $t = -2.3$)---it
pays off when dealer credit deteriorates---and positively on the
default premium ($+0.11$, $t = 3.2$) and the industrial bond return
($+0.22$, $t = 3.1$), the natural exposures of a leveraged
long-credit position, with smaller loadings on high yield,
emerging-market debt, the currency trend-following factor, and the
curvature trade. The iTraxx set is
anchored by the traded intermediary capital factor of
\citet{he2017intermediary} ($-0.03$, $t = -4.1$)---the strategy pays
off when dealer equity underperforms---alongside
betting-against-beta ($-0.04$, $t = -4.3$), global and fixed-income
momentum ($-0.05$ and $-0.04$, $t = -3.1$ and $-3.9$), the one-month
variance-swap return ($-0.86$, $t = -2.5$), value, and
investment-grade credit risk through the CDX index. The three selected
sets are pairwise
disjoint: no factor appears in more than one strategy, so the alpha
that survives is not compensation for a single common exposure, and
its origin is strategy-specific.

\paragraph{Out-of-Sample Alpha.}

Panel~B of Table~\ref{tab:pca_oos} applies the recursive design of
Section~\ref{sec:pca} to the selected models: the best-subset factor
set is frozen, only the hedge loadings are re-estimated on the
past-only window, and the realised hedged return is tested with the
HAC correction---the real-time abnormal return of the selected model,
with no per-window re-selection. The out-of-sample alpha is positive
and significant at the $1\%$ level for all three strategies under
both schemes: $+2.05\%$ p.a.\ for BTP~Italia, $+3.90\%$ for
CDS--Bond Basis, and $+2.50\%$ for iTraxx Skew on the expanding
window, against in-sample estimates of $+2.2$, $+4.8$, and $+2.9\%$
in Table~\ref{tab:aen_stable_ols}.

\paragraph{Selector Robustness.}

Table~\ref{tab:selector_robustness} shows that neither the $\ell_0$
machinery nor the post-selection inference drives the result. Five
alternative selectors span the methodological space: a
best subset constrained to one factor per risk category, the adaptive
elastic net of \citet{zou2009adaptive}, the relaxed lasso of
\citet{meinshausen2007relaxed}, a nonlinear machine-learning screen
(random forest permutation importance, \citealp{breiman2001random}),
and model-X knockoffs \citep{candes2018panning} with the
false discovery rate controlled at $10\%$. For every selector the
table reports the honest split-sample alpha---factors selected on the
first half of the sample, alpha estimated on the second---and the
frozen-set out-of-sample alpha. Every estimate for the CDS--Bond basis
and the iTraxx skew is significant at the $1\%$ level and every
estimate for BTP~Italia at the $5\%$ level, with the honest estimates
spanning $+1.5$ to $+2.0\%$ p.a.\ for BTP~Italia, $+2.5$ to $+3.4\%$ for the CDS--Bond basis, and
$+2.1$ to $+2.5\%$ for iTraxx Skew; for the primary selector the
out-of-sample column coincides with Panel~B of
Table~\ref{tab:pca_oos}. The knockoffs row is the sharpest statement
of the two-sided message: at a controlled false discovery rate, no
individual candidate carries enough signal to be selected at all, and
the resulting no-controls alpha is the raw strategy alpha---so the
data neither identify a stable factor set nor produce one that
absorbs the premium. The identity of the selected factors is
unstable across selectors and sample halves, but the exposures that
recur are the consensus ones---the fixed-income defensive factor for
BTP~Italia, emerging-market debt and the industrial bond return for
the CDS--Bond basis, the intermediary and defensive-style block for
iTraxx Skew---and the double-selection control set is a strict subset
of exactly these, empty for the iTraxx skew, where no candidate clears
its threshold (Appendix~\ref{app:aen_selection_matrix}). The out-of-sample alpha of
the best-subset hedge clears in all three strategies, under both window
designs, the $t>3.0$ hurdle that \citet{harvey2016and} propose for
factors discovered in multiple-testing environments
(Table~\ref{tab:pca_oos}, Panel~B).


% ── III.D CONDITIONAL ALPHA ───────────────────────────────────────────────────

\subsection{Conditional Alpha}
\label{sec:conditional_alpha}

Table~\ref{tab:alpha_regime_pca_aen} decomposes alpha across the three
stress regimes on the iTraxx Main 5-year level---LOW, MEDIUM, and
HIGH, with the same 60 and 100 basis-point cutoffs used throughout the
paper---under both the PCA and the best-subset controls. In every
cell, alpha rises monotonically from the LOW to the HIGH regime. The
CDS--Bond basis shows the sharpest gradient---from about $+2.8\%$ in
LOW to $+7.6\%$ in HIGH under either set of controls---and BTP~Italia the
familiar pattern of an alpha concentrated away from calm markets:
under the PCA controls it rises from $+0.4\%$ (insignificant) in LOW
to $+3.3\%$ in HIGH; under the best-subset controls the whole
profile shifts up, from $+1.3\%$ to $+3.8\%$. For iTraxx Skew
the alpha is significant in every regime under both specifications,
rising from $+1.0\%$ to $+3.5\%$ under the PCA controls, with a
higher and flatter profile ($+2.5\%$ to $+3.4\%$) under the
best-subset controls.

The \citet{christopherson1998conditioning} conditional model of
Section~\ref{sec:benchmarks}---the full specification of
equation~(\ref{eq:cfg}), with the lagged standardised stress state
conditioning both the alpha and the loadings, and $\alpha_1$ tested by
the moving-block bootstrap---is re-estimated on each factor layer in
Appendices~\ref{app:pca_conditional} and~\ref{app:aen_conditional}.
Under the PCA controls, $\alpha_1$ is significant at the $1\%$ level
for all three strategies. Under the best-subset controls, it remains
large and significant for the CDS--Bond basis ($+2.55\%$ per standard
deviation, $t = 4.31$), while for BTP~Italia and iTraxx Skew the
continuous slope is absorbed once the loadings themselves are allowed
to vary with stress---the regime gradient of
Table~\ref{tab:alpha_regime_pca_aen} survives, but it is carried by
the conditional loadings rather than by a conditional intercept.
Sub-period stability of the alpha estimates is confirmed in
Appendix~\ref{app:pca_subperiod}, and rolling 36-month alpha estimates
with regime shading are plotted in Appendix~\ref{app:aen_rolling}.


% ── III.E ALPHA ACROSS LAYERS ─────────────────────────────────────────────────

\subsection{Alpha Across Layers}
\label{sec:factor_discussion}

Table~\ref{tab:aen_method_comparison} consolidates the answer to the
question this section and the previous one pose. Each cell reports the
alpha and adjusted $\bar{R}^2$ from one layer of controls: the three
benchmark frameworks of Section~\ref{sec:benchmarks}, the principal
components, and the best-subset selection. Moving down the table, the
explanatory power rises---from $\bar{R}^2$ between $0\%$ and $15\%$
under the benchmarks to $21$--$31\%$ under the best-subset
models---while the alpha barely moves: $+1.4$ to $+2.2\%$ p.a.\ for
BTP~Italia, $+4.0$ to $+4.8\%$ for the CDS--Bond basis, and $+1.8$ to
$+2.9\%$ for iTraxx Skew, significant at the $1\%$ level in every
cell. A supervised search over $75$ tradable candidates, free to pick
whichever eight factors best explain each strategy, absorbs no more of
the alpha than the fixed benchmarks do.

Two conclusions follow. First, the excess returns documented in
Section~\ref{sec:performance} are not compensation for systematic risk
exposures, whether those exposures are specified ex ante, extracted as
latent components, or selected by the data. Second, the alpha that
survives is state-dependent, rising when intermediary funding stress
is elevated, under every layer of controls. The natural question is
why such a premium persists---and whether the three mispricings that
generate it are connected through the balance sheets of the
intermediaries that trade them. Section~\ref{sec:interdependencies}
takes up that question.